\SourceLeaf{427}%
das Element $\delta\varrho$, so wird:
\[
\delta x=X\delta\varrho,\qquad\delta y=Y\delta\varrho,\qquad\delta z=Z\delta\varrho
\]
und
\[
\delta V=P\delta x+Q\delta y+R\delta z=\mu\delta\varrho.
\]
Man sieht also, wie das Zeichen von $\mu$ von der Zeichenänderung des Werthes von $V$ beim Durchgang von der untern auf die obere Seite abhängt.

\begin{center}9.\end{center}
\noindent\hspace*{1.2em}Man schneide die krumme Fläche mit einer Ebene in dem Punkte, auf den sich unsere Bezeichnungen beziehen; es entsteht so eine plane krumme Linie, von welcher $ds\mathord{\hbox{\begin{tikzpicture}[x=.1pt,y=-.1pt,baseline=-7pt]
\path[use as bounding box] (0,0) rectangle (17,101);
\path[fill=black] (14,69) -- (12,69) -- (12,70) -- (9,72) -- (9,74) -- (3,80) -- (2,82) -- (2,85) -- (0,87) -- (0,88) -- (2,90) -- (2,96) -- (3,97) -- (3,99) -- (5,100) -- (11,100) -- (14,97) -- (14,88) -- (13,87) -- (13,84) -- (14,83) -- (14,78) -- (16,76) -- (16,75) -- (14,73) -- cycle (6,3) -- (5,4) -- (5,12) -- (6,14) -- (9,14) -- (14,9) -- (14,6) -- (11,3) -- cycle;
\end{tikzpicture}
}}$ ein Element sei, in Beziehung auf welches wir die obigen Bezeichnungen beibehalten. Als obere Seite der Ebene betrachten wir die, auf welcher die Normale auf die krumme Fläche liegt; auf die Ebene errichten wir eine Normale, deren Richtung durch den Punkt $\mathfrak{L}$ der Hülfskugel bezeichnet werde. Beim Fortschreiten in der krummen Linie werden also $\lambda$ und $L$ ihre Stelle verändern, während $\mathfrak{L}$ constant bleibt und immer $\lambda L$ und $\lambda\mathfrak{L}=90^{0}$ sind; $\lambda$ beschreibt daher den grössten Kreis, dessen einer Pol $\mathfrak{L}$ ist. Das Element dieses grössten Kreises wird
\[
=\frac{ds}{r}
\]
sein, wenn $r$ den Krümmungshalbmesser der Curve bedeutet; und wird wiederum die Richtung eben dieses Elements auf der Hülfskugel mit $\lambda'$ bezeichnet, so wird offenbar $\lambda'$ in eben dem grössten Kreise liegen, und sowohl von $\lambda$ als von $\mathfrak{L}$ um $90^{0}$ abstehen. Setzt man nun noch
\[
\cos\lambda'(1)=\xi',\qquad\cos\lambda'(2)=\eta',\qquad\cos\lambda'(3)=\zeta',
\]
so wird
\[
d\xi=\xi'\frac{ds}{r},\qquad d\eta=\eta'\frac{ds}{r},\qquad d\zeta=\zeta'\frac{ds}{r}
\]
sein, da in der That $\xi,\eta,\zeta$ nichts anderes sind, als die Coordinaten des Punktes $\lambda$ gegen den Mittelpunkt der Kugel.
